\paragraph{Articulation estimation.}
We initialize articulation from the moving part's 3D point trajectories
and refine it through differentiable rendering. We first re-anchor the
trajectories to the mesh reference frame and estimate rigid transformations
at training keyframes using trimmed Kabsch alignment. We robustly fit
prismatic and revolute joint hypotheses to these transformations:
\begin{align}
    \mathcal{T}_{\mathrm{P}}(\mathbf{x};q_t)
    &= \mathbf{x}+q_t\mathbf{a}, \\
    \mathcal{T}_{\mathrm{R}}(\mathbf{x};q_t)
    &= \mathbf{c}+\exp\!\left(q_t[\mathbf{a}]_\times\right)
       (\mathbf{x}-\mathbf{c}),
\end{align}
where $\mathbf{a}$ is the unit joint axis, $\mathbf{c}$ is a point on
the revolute axis, and $q_t$ is the articulation state. After aligning
the reference meshes using the static-part masks, we refine each
initialized joint and its states by minimizing silhouette, boundary,
and DINOv2 feature discrepancies between rendered and observed images,
with temporal smoothness regularization and a shared object pose and
scale correction. The 3D tracks are used only for initialization.
We select the joint type with the lower median image loss on validation
frames, optimizing only their states during evaluation. Finally, we fix
the selected joint model, recover states for all frames, and temporally
smooth the resulting sequence.
